\documentclass[11pt]{article}
\usepackage[margin=.92in]{geometry}
\usepackage{amsmath,amssymb,booktabs,hyperref}
\hypersetup{colorlinks=true,linkcolor=black,urlcolor=blue}
\title{A missing prime in the finite example for OEIS A300984}
\author{Blake Stone's mathematics project\\\small Developed with OpenAI Codex}
\date{November 27, 2026}
\begin{document}
\maketitle
For a positive integer $n$, write $S(n)$ and $T(n)$ for the sums of its
squarefree and nonsquarefree divisors, respectively. Michel Lagneau's A300984
consists of the integers for which both sums are squarefree. The 2018 entry
lists the odd prime bases occurring in its members of size at most $10^7$ and
gives the proportion belonging to A082539 as $84\%$. That list omits $937$.
The corrected count is $39$ of $47$, or $82.98\%$ to two decimal places.

\section*{An explicit missing base}
Take $n=4\cdot937^2=3{,}511{,}876$. The squarefree divisors are $1,2,937,1874$.
For the remaining divisors use multiplicativity of the ordinary divisor sum:
\[
 S(n)=2814=2\cdot3\cdot7\cdot67,\qquad
 T(n)=7(1+937+937^2)-2814=6{,}149{,}535=3\cdot5\cdot7\cdot58{,}567.
\]
All the displayed factors are distinct primes in their respective products
(primality can be checked by trial division through the square root).
A prime $p$ belongs to A082539 exactly when no smaller prime $r$ has
$r+1\mid p+1$. Since $13+1\mid937+1$, the omitted base contributes to the
denominator only. A separate divisor-by-divisor enumeration verifies that the
entire printed set together with $937$ is exactly the set at this cutoff.

\section*{Exact counts at larger decimal cutoffs}
The prime-support theorem in the accompanying project shows that every member
of A300984 has the form $2^i p^j$ with $i,j\ge2$, $p\equiv1\pmod{12}$,
and $3(p+1)$ squarefree. For candidates of this form,
\[
 S(2^i p^j)=3(p+1),\qquad
 T(2^i p^j)=(2^{i+1}-1)(1+p+\cdots+p^j)-3(p+1).
\]
These two equalities give a finite exhaustive test. For cutoff $X$, sieve primes
$p\le\lfloor\sqrt{X/4}\rfloor$. For each, find the smallest candidate for which
$T$ is squarefree by enumerating both exponents within $2^i p^j\le X$.
After finding a witness, larger candidates for the same base can be dropped.
Mark failures of A082539 by striking the larger multiples of $r+1$ for all
primes $r$ in the sieve. Each base is counted once at each cutoff that contains its first witness.

\begin{center}\begin{tabular}{r@{\quad}rr|r@{\quad}rr}
\toprule
$m$ & Bases at $10^m$ & In A082539 & $m$ & Bases at $10^m$ & In A082539\\
\midrule
7 & 47 & 39 & 12 & 8,880 & 5,670\\
8 & 132 & 98 & 13 & 25,639 & 15,783\\
9 & 382 & 269 & 14 & 74,802 & 44,947\\
10 & 1,084 & 750 & 15 & 218,702 & 128,324\\
11 & 3,087 & 2,053 & 16 & 644,008 & 369,981\\
\bottomrule\end{tabular}\end{center}
Thus the conjectured strict majority holds at the nine cutoffs $10^8$ through
$10^{16}$. At the largest cutoff its share is about $57.45\%$. The assertion
for all integers $m>7$ remains open.

\paragraph{Reproduction.}
\texttt{majority.cpp} generates the full table using integer divisibility
and square roots corrected by integer comparisons. To test squarefreeness it
removes trial primes once up to the cube root of the current cofactor,
rejecting any repeated factor. On termination the residual has at most two
prime factors with multiplicity, so it suffices to test for a perfect square
greater than one. The independent \texttt{check\_majority.py} enumerates each
candidate's divisors, factors their two sums using SymPy, and tests the proper
divisors of $p+1$. It checks the witness, the complete printed set, and the
counts through $10^9$. The earlier shape proof is in \texttt{Proof.lean}.
The source entries are \href{https://oeis.org/A300984}{A300984} and
\href{https://oeis.org/A082539}{A082539}.
\end{document}
